\documentclass[a4paper,11pt]{jsarticle}
\usepackage[dvipdfmx]{graphicx}
\usepackage{geometry}
\usepackage{amsmath}
\usepackage{booktabs}
\usepackage{longtable}
\usepackage{siunitx}
\geometry{margin=22mm}

\title{p2MEG 最終解析における ACC 時間構造のデータ駆動モデル化\\
step4f\_run1to20\_eg\_doubleonly\_allpatterns\_500ns.txt の検討}
\author{p2MEG 解析レポート}
\date{2026年3月6日}

\begin{document}
\maketitle

\begin{abstract}
最終解析用再構成データ
\texttt{data/finaldata/step4f\_run1to20\_eg\_doubleonly\_allpatterns\_500ns.txt}
を用いて、ACC（accidental）事象の時間構造をデータ駆動で評価した。
当初想定していた
\(
p_{\mathrm{acc}}(x)=p_4(E_e,E_\gamma,\phi_e,\phi_\gamma)\times p_t(t)
\)
かつ \(p_t(t)=\) 一様
という仮定は、実データの時間分布から成立しない。
そこで、まず \(E_\gamma\) ごとの時間分布 \(p_t(t|E_\gamma)\) を broad な対称関数でフィットし、
各 \(E_\gamma\) bin について
\(
R_i = N_i(\mathrm{AW}) / N_i(\mathrm{TSB})
\)
を求めた。
次に、この \(R_i\) を timing sideband (TSB) の実測事象数に掛けて
解析窓 (AW) 内の ACC 数を予測した。
その結果、現在の解析窓
\(
20\le E_e\le 80\;\mathrm{MeV},
20\le E_\gamma\le 80\;\mathrm{MeV},
-50\le t\le 50\;\mathrm{ns},
1.7\le \theta_{eg}\le \pi
\)
では、観測事象数 \(18\) に対し
\(
N_{\mathrm{ACC,pred}} = 17.85 \pm 2.82
\)
となり、統計誤差と簡易モデル不定性の範囲で
現状の解析窓はほぼ全て ACC で説明できることが分かった。
\end{abstract}

\section{目的}
本解析の目的は、最終 5D データだけを用いて ACC 時間構造をできる限り定量化し、
最終解析に使う background 見積りの現実的な基礎を与えることである。

今回利用できた情報は最終再構成後の
\[
(E_e,\; E_\gamma,\; t,\; \phi_e,\; \phi_\gamma)
\]
のみであり、再構成アルゴリズムや event-by-event の共通時間原点は不明であった。
したがって、raw レベルの hit mixing や厳密な event mixing は行えない。
その代わりに、「全データの大部分は ACC である」という仮定の下で、
時間分布そのものを broad な ACC 構造として扱い、
\(E_\gamma\) ごとの依存性を取り込む方針を採用した。

\section{問題設定}
全データのヒストグラムから、次の 2 点が問題であることが分かっていた。

\begin{enumerate}
  \item \(t\) 分布は一様ではなく、\(t=0\) 付近に幅 \(\mathcal{O}(100\,\mathrm{ns})\) の broad peak を持つ。
  \item \(N(\mathrm{AW})/N(\mathrm{TSB})\) は \(E_\gamma\) に依存する。
\end{enumerate}

このため、従来の
\[
N_{\mathrm{ACC}}(\mathrm{AW})
=
\frac{W_{\mathrm{AW}}}{W_{\mathrm{TSB}}} N_{\mathrm{TSB}}
\]
という単純スケーリングは不適切である。

\section{入力データと選別}
\subsection{入力ファイル}
解析対象は
\[
\texttt{data/finaldata/step4f\_run1to20\_eg\_doubleonly\_allpatterns\_500ns.txt}
\]
である。
5 列は
\[
E_e,\;E_\gamma,\;t,\;\phi_e,\;\phi_\gamma
\]
であり、単位はそれぞれ MeV, MeV, ns, rad, rad とした。

\subsection{角度と解析窓}
角度は
\[
\theta_{eg} = |\phi_e-\phi_\gamma|
\]
で定義した。
今回の解析で用いた最終解析窓はコード設定
\texttt{include/p2meg/AnalysisWindow.h}
に合わせて
\begin{align}
20 &\le E_e \le 80 \;\mathrm{MeV},\\
20 &\le E_\gamma \le 80 \;\mathrm{MeV},\\
-50 &\le t \le 50 \;\mathrm{ns},\\
1.7 &\le \theta_{eg} \le \pi
\end{align}
とした。
TSB は
\[
t\in[-500,500]\;\mathrm{ns}
\quad\text{かつ}\quad
t\notin[-50,50]\;\mathrm{ns}
\]
で定義した。

\section{ACC 時間構造のモデル化}
\subsection{基本方針}
再構成アルゴリズムが不明なため、prompt 核の厳密形ではなく、
全データの時間分布を broad な ACC 構造として smooth に近似した。
特に \(E_\gamma\) 依存が強いため、\(E_\gamma\) を
\[
[0,10],\ [10,20],\ [20,30],\ [30,40],\ [40,50],\ [50,80]\ \mathrm{MeV}
\]
の bin に分けて時間分布を作成した。

\subsection{フィット関数}
各 \(E_\gamma\) bin に対して、\(|t|<20\;\mathrm{ns}\) を blind core とし、
その外側の分布を次の対称関数でフィットした。
\begin{equation}
f_i(t)=A_i\exp\left(-\frac{t^2}{2\sigma_i^2}\right)+C_i.
\end{equation}
ここで \(A_i\) は振幅、\(\sigma_i\) は幅、\(C_i\) は pedestal である。

blind core の縁だけを拾う不自然に鋭い peak を避けるため、
\[
\sigma \ge 20\ \mathrm{ns}
\]
を下限として課した。
これは signal/RMD の prompt 核を記述するためではなく、
あくまでビーム起源の broad ACC 構造だけを見るための制約である。

\subsection{フィットの目視確認}
フィット結果は
\texttt{doc/finalanalysis/fit\_acc\_time\_by\_eg\_step4f\_run1to20\_eg\_doubleonly\_allpatterns\_500ns.pdf}
に保存した。
代表として \(E_\gamma=0\text{--}10\), \(10\text{--}20\), \(20\text{--}30\), \(30\text{--}40\) MeV の結果を
図\ref{fig:fitpages} に示す。

\begin{figure}[p]
  \centering
  \includegraphics[width=0.9\linewidth,page=2]{fit_acc_time_by_eg_step4f_run1to20_eg_doubleonly_allpatterns_500ns.pdf}
  \vspace{3mm}
  \includegraphics[width=0.9\linewidth,page=3]{fit_acc_time_by_eg_step4f_run1to20_eg_doubleonly_allpatterns_500ns.pdf}
  \caption{\(E_\gamma=0\text{--}10\) MeV および \(10\text{--}20\) MeV の時間分布フィット。灰色帯は blind core (\(|t|<20\) ns)。}
  \label{fig:fitpages}
\end{figure}

\begin{figure}[p]
  \centering
  \includegraphics[width=0.9\linewidth,page=4]{fit_acc_time_by_eg_step4f_run1to20_eg_doubleonly_allpatterns_500ns.pdf}
  \vspace{3mm}
  \includegraphics[width=0.9\linewidth,page=5]{fit_acc_time_by_eg_step4f_run1to20_eg_doubleonly_allpatterns_500ns.pdf}
  \caption{\(E_\gamma=20\text{--}30\) MeV および \(30\text{--}40\) MeV の時間分布フィット。}
  \label{fig:fitpages2}
\end{figure}

目視では、
\begin{itemize}
  \item \(0\text{--}10\) MeV では broad peak を明確に再現している。
  \item \(10\text{--}20\) MeV と \(20\text{--}30\) MeV でも不自然な狭い中心 peak は抑えられた。
  \item \(30\text{--}40\) MeV は統計が薄く、ほぼ flat に近いが破綻はしていない。
  \item \(40\) MeV 以上は統計が小さすぎて、定量的にはほぼ参考値である。
\end{itemize}
と判断した。

\section{\(R_i = N_i(\mathrm{AW}) / N_i(\mathrm{TSB})\) の導出}
各 \(E_\gamma\) bin について、フィット関数から
\begin{equation}
R_i^{\mathrm{fit}} =
\frac{
\int_{-50}^{50} f_i(t)\,dt
}{
\int_{-500}^{-50} f_i(t)\,dt
+
\int_{50}^{500} f_i(t)\,dt
}
\end{equation}
を定義した。
比較のため、生データからの
\[
R_i^{\mathrm{data}} = N_i(\mathrm{AW}) / N_i(\mathrm{TSB})
\]
も併記した。

結果を表\ref{tab:rvalues} に示す。

\begin{table}[htbp]
\centering
\caption{\(E_\gamma\) bin ごとの \(R_i\) の比較}
\label{tab:rvalues}
\begin{tabular}{cccccc}
\toprule
\(E_\gamma\) [MeV] & entries & \(N_{\mathrm{AW}}\) & \(N_{\mathrm{TSB}}\) & \(R_i^{\mathrm{data}}\) & \(R_i^{\mathrm{fit}}\) \\
\midrule
0--10  & 331 & 111 & 220 & 0.5045 & 0.5501 \\
10--20 &  92 &  21 &  71 & 0.2958 & 0.2079 \\
20--30 &  79 &  16 &  63 & 0.2540 & 0.2384 \\
30--40 &  23 &   1 &  22 & 0.0455 & 0.1135 \\
40--50 &   2 &   1 &   1 & 1.0000 & 0.1136 \\
50--80 &   2 &   0 &   2 & 0.0000 & 0.1114 \\
\bottomrule
\end{tabular}
\end{table}

最終解析窓では \(E_\gamma \ge 20\) MeV しか使わないため、
実際に寄与するのは主に
\(20\text{--}30\) MeV, \(30\text{--}40\) MeV, \(40\text{--}50\) MeV, \(50\text{--}80\) MeV
の4 bin である。

\section{TSB からの ACC 数予測}
\subsection{定義}
各 \(E_\gamma\) bin について、
\begin{equation}
N_{\mathrm{ACC,pred},i}(\mathrm{AW})
=
R_i^{\mathrm{fit}} N_i(\mathrm{TSB})
\end{equation}
で ACC 事象数を予測した。
総和は
\begin{equation}
N_{\mathrm{ACC,pred}}(\mathrm{AW})
=
\sum_i N_{\mathrm{ACC,pred},i}(\mathrm{AW})
\end{equation}
とした。

\subsection{誤差の付け方}
現在のデータだけで厳密な系統誤差は定義できないため、
各 bin で次の 2 成分を足し合わせた。
\begin{enumerate}
  \item TSB 計数の Poisson 統計誤差
  \[
  \sigma_{\mathrm{stat},i}
  =
  R_i^{\mathrm{fit}}\sqrt{N_i(\mathrm{TSB})}
  \]
  \item smooth fit と実測比の差を使った簡易モデル不定性
  \[
  \sigma_{\mathrm{model},i}
  =
  N_i(\mathrm{TSB})
  \left|R_i^{\mathrm{fit}}-R_i^{\mathrm{data}}\right|
  \]
\end{enumerate}
したがって、
\begin{equation}
\sigma_i^2
=
\sigma_{\mathrm{stat},i}^2+\sigma_{\mathrm{model},i}^2
\end{equation}
とした。
全体誤差は bin 間独立近似の下で二乗和平方根を取った。

\subsection{結果}
結果を図\ref{fig:predacc} に示す。

\begin{figure}[htbp]
  \centering
  \includegraphics[width=0.95\linewidth,page=1]{predict_acc_from_tsb_by_eg_step4f_run1to20_eg_doubleonly_allpatterns_500ns.pdf}
  \vspace{3mm}
  \includegraphics[width=0.95\linewidth,page=2]{predict_acc_from_tsb_by_eg_step4f_run1to20_eg_doubleonly_allpatterns_500ns.pdf}
  \caption{TSB からの ACC 数予測。上: 全体 summary。下: \(E_\gamma\) bin ごとの AW 観測数と ACC 予測数。}
  \label{fig:predacc}
\end{figure}

数値は次の通りである。
\begin{align}
N_{\mathrm{AW,obs}} &= 18,\\
N_{\mathrm{ACC,pred}} &= 17.85 \pm 2.82,\\
N_{\mathrm{AW,obs}} - N_{\mathrm{ACC,pred}} &= 0.15.
\end{align}

bin ごとの主な内訳は
\begin{itemize}
  \item \(20\text{--}30\) MeV: \(16\) observed に対し \(15.02 \pm 2.13\) predicted
  \item \(30\text{--}40\) MeV: \(1\) observed に対し \(2.50 \pm 1.59\) predicted
  \item \(40\text{--}50\) MeV: \(1\) observed に対し \(0.11 \pm 0.89\) predicted
  \item \(50\text{--}80\) MeV: \(0\) observed に対し \(0.22 \pm 0.27\) predicted
\end{itemize}
である。

\section{考察}
\subsection{今回の結論}
現在の最終解析窓では、観測された AW 事象数は
データ駆動で求めた ACC 予測と整合している。
少なくとも event count のレベルでは、
「現時点で見えている事象はほぼ全て ACC で説明できる」
というのが最も保守的で妥当な結論である。

\subsection{この結果の意味}
今回の結果は「信号が全く無いことの証明」ではない。
しかし、
\[
N_{\mathrm{AW,obs}} - N_{\mathrm{ACC,pred}}
= 0.15
\]
であり、その大きさは今回付けた誤差 \(\pm 2.82\) より十分小さい。
したがって、現在の解析窓で統計的に有意な excess を主張することはできない。

\subsection{限界}
本解析には次の制約がある。
\begin{enumerate}
  \item 再構成アルゴリズムと event-level 時間原点が不明であり、
  raw に戻った厳密な hit mixing ができない。
  \item 誤差評価のうち \(\sigma_{\mathrm{model}}\) は簡易的な proxy であり、
  厳密な系統誤差ではない。
  \item \(E_\gamma \ge 40\) MeV の bin は統計が少なく、\(R_i\) は不安定である。
\end{enumerate}

\subsection{次にやるべきこと}
今後の優先順位は次の通りである。
\begin{enumerate}
  \item \(E_\gamma\) だけでなく \(\theta_{eg}\) 依存も含めた
  \(R(E_\gamma,\theta)\) を調べる。
  \item Eg 依存スケールで TSB を重み付きにした差し引きヒストグラムを作り、
  residual の形状を確認する。
  \item もし再構成情報が将来得られるなら、raw あるいは singles レベルで
  ACC 時間構造を作り直す。
\end{enumerate}

\section{まとめ}
最終 5D データのみを用いて ACC 時間構造をデータ駆動で評価した。
blind core を除いた \(E_\gamma\) ごとの時間分布を smooth な対称関数でフィットし、
それに基づいて \(R_i = N_i(\mathrm{AW})/N_i(\mathrm{TSB})\) を求めた。
この \(R_i\) を使って TSB から AW 内の ACC 数を予測した結果、
\[
N_{\mathrm{ACC,pred}} = 17.85 \pm 2.82
\]
となり、観測 \(18\) 事象と整合した。
従って、現時点では現在の最終解析窓の事象は
ACC でほぼ説明できると結論するのが妥当である。

\end{document}
